% ==================================================
% 1. Abstract and introduction
%
% Deliberately no numbered subsections: cleanthesis section headings are tall, and an
% introduction this short does not need a table of contents of its own. Paragraph breaks
% carry the structure instead.
% ==================================================
\chapter{Abstract and Introduction}
\label{sec:introduction}

\section*{Abstract}

We add three recent Shapley-value approximators to \shapiq{} --- \leverageshap{}
\citep{Witter.2025a}, \polyshap{} \citep{Fumagalli.2026a} and \oddshap{}
\citep{Fumagalli.2026} --- reproduce each paper from the resulting implementation, and
benchmark all three against every approximator \shapiq{} already provides on one shared
testbed. Each paper's headline claim holds, and holds on a value function none of the three
papers uses: \oddshap{} leads on the high-dimensional games once the budget passes the
threshold its paper predicts, \leverageshap{} beats optimised KernelSHAP at every dimension,
and \polyshap{} is strongest at low dimension and degrades at high, for the reason its own
paper gives. The comparison also shows that no single ordering of the three exists. Averaged
over all games \leverageshap{} ranks best; restricted to the games above sixty players
\oddshap{} does, by a wide margin. Which of those two numbers is the answer depends on the
dimension a reader cares about, and reporting either alone would misrepresent the other.

% --------------------------------------------------
% The introduction runs as continuous prose. Roughly one paragraph per block.
% --------------------------------------------------

\paragraph{What this report is.}
The Shapley value assigns each feature its average marginal contribution over all $2^d$
coalitions, which is why every practical use of it rests on an approximator. This report covers
three recent ones. We implemented each in \shapiq{} \citep{Muschalik.2024b}, reproduced the
experiments of the paper that introduced it, and then compared all three against the
approximators \shapiq{} already had. Chapter~\ref{sec:benchmark} is the comparison that puts
every method on the same footing; Chapters~\ref{sec:oddshap} to~\ref{sec:polyshap} take each
approximator on its own terms, first summarising its paper, then reproducing it.

\paragraph{What we did.}
Three approximators were implemented and submitted upstream. \oddshap{} is merged
(PR~\#522, with a follow-up in PR~\#560); \leverageshap{} (PR~\#524) and \polyshap{}
(PR~\#559) are open at the time of writing. Each has a reproduction of its paper, driven by
committed result tables so that every figure re-renders offline in minutes, without a GPU and
without re-running the experiments. On top of those, the cross-method benchmark of
Chapter~\ref{sec:benchmark} runs every approximator registered in \shapiq{} over a shared set
of eleven games, twelve budgets and ten explained predictions --- 19{,}800 cells in total.
Two defects in \shapiq{} surfaced while the benchmark was being built and are fixed alongside
this work: the ViT value function moved its model to the GPU but left its normalisation
parameters on the CPU, so it could not run there at all; and the placeholder classes that stand
in for the optional \texttt{sparse} and \texttt{shapleig} extras were abstract, so a missing
dependency was reported as a broken approximator rather than as a missing dependency.

\paragraph{How to read it, and one distinction that runs through it.}
The report answers two different questions, and they need different setups. A
\emph{reproduction} asks whether a paper's claim survives re-implementation; it has to use that
paper's own value function, or it tests nothing. The three papers do not share one --- \oddshap{}
alone trains XGBoost classifiers and defines its tabular value function by interventional
perturbation over fifty background instances. The \emph{benchmark} asks how the methods compare
with each other and with what \shapiq{} already had; it has to use a single value function for
every method, which then belongs to no paper in particular. The consequence is worth stating
before any table is read: the absolute errors in Chapter~\ref{sec:benchmark} and those in
Chapters~\ref{sec:oddshap} to~\ref{sec:polyshap} are different quantities, measured against
different targets, and comparing them across chapters is meaningless. Within each chapter they
are directly comparable, which is what each chapter is for. That the papers' claims survive the
move to a value function none of them chose is the stronger result, and it is only visible
because the two settings are kept apart.
